\section{Sonuçlar}
\label{sec:results}

Sonuçları L-S-Q boru hattındaki sıraya göre raporluyoruz
(Bölüm~\ref{sec:results:lsq}--\ref{sec:results:quant}); ardından
gömülü hedeflerdeki çalışma davranışına geçiyoruz
(Bölüm~\ref{sec:results:crossplatform}--\ref{sec:results:streaming}).
Tüm deneyler Bölüm~\ref{sec:setup}'daki protokolü kullanır.

\subsection{L-S-Q Boru Hattı Doğruluğu}
\label{sec:results:lsq}

Tablo~\ref{tab:lsq}, her sıkıştırma aşamasının birikimli etkisini
özetler; Tablo~\ref{tab:multiseed} tohum bazlı tam dökümü verir.
Sınıflandırıcı başlığıyla birlikte tam ranklı FastGRNN
hücresi (Bölüm~\ref{sec:method:cell}) toplam 440 parametreye
sahiptir ve test kümesinde F1~$=$~0.912 değerine ulaşır. Düşük
ranklı ayrıştırma ($r_w{=}2$, $r_u{=}8$), yinelemeli matris
parametre sayısını $\mathbf{U} = \mathbf{U}_1 \mathbf{U}_2^\top$
ayrıştırması yoluyla azaltır (beş tohum üzerinde $0.879 \pm 0.052$).
$s{=}0.5$ düzeyindeki seyrekleştirme sayıyı 283 sıfır olmayan
parametreye indirir ($0.856 \pm 0.099$ ortalama). Tensör-başına Q15
nicelendirme, beş tohumda $\mathbf{0.853 \pm 0.107}$ makro F1
verir; gömülü seed-0 modeli bare-metal \msp{} üzerinde
\textbf{0.9176} F1'e ulaşır (C motoru, PyTorch ile \%100 pencere
uyuşması). Büyük standart sapma, tüm aşamalarda tutarlı biçimde
düşük kalan tohum~1 yakınsama aykırısından kaynaklanır
(Tablo~\ref{tab:multiseed}, Şekil~\ref{fig:lowrank}).
Tohum~1, Q15 F1'i (0.663) MLP taban çizgisinin (0.847) altına
düşen tek tohumdur; kalan dört tohumun tamamı 0.89'un üzerinde
kalarak MLP'yi her sıkıştırma aşamasında belirgin biçimde geçmektedir.
Şekil~\ref{fig:saturation}, sonraki tüm aşamalara temel oluşturan
$H{=}16$ eğitim çalışmasını gösterir.

\begin{figure}[t]
\centering
\includegraphics[width=\linewidth]{figures/saturation.pdf}
\caption{$H{=}16$ gizli boyutu için eğitim yörüngesi.
Test F1 epoch~100 civarında doyar; en iyi tohum kontrol noktası
(epoch~113'te F1~$=$~0.912), sonraki tüm sıkıştırma aşamaları için
kullanılan temel modeldir.}
\label{fig:saturation}
\end{figure}

\begin{table}[t]
\centering
\caption{\hapt{} üzerinde ardışık L-S-Q boru hattı. Her satır,
belirtilen dönüşümü birikimli olarak uygular. \emph{Sıfır-olmayan},
hücre ve yoğun sınıflandırıcı başlığı dahil tüm sıfır olmayan
parametreleri sayar; başlık her aşamada 102 yoğun parametre katar.
Ortalama F1 beş tohum üzerindendir; son Q15 satırı gömülü hedefe
aktarılan seed-0 modeli üzerindedir. Boyut ilk üç satırda FP32,
son satırda Q15'tir.}
\label{tab:lsq}
\setlength{\tabcolsep}{4pt}
\begin{tabular}{@{}lrrr@{}}
\toprule
Aşama & F1 & Sıfır-olmayan & Boyut \\
\midrule
FastGRNN tam rank ($H{=}16$)        & 0.912          & 440 & 1.8~KB \\
+ Düşük rank ($r_w{=}2$, $r_u{=}8$) & 0.879          & 430 & 1.7~KB \\
+ Seyreklik ($s{=}0.5$)             & 0.856          & 283 & 1.1~KB \\
+ Q15 (ağırlık + kalib.\ akt.)      & \textbf{0.918} & 283 & \textbf{566~B} \\
\bottomrule
\end{tabular}
\end{table}

\begin{table}[t]
\centering
\caption{\hapt{} üzerinde tohum bazlı L-S-Q boru hattı (3{,}399 test
penceresi). \emph{Q15/LUT F1}: PyTorch benzetimi, Q15 ağırlık +
kalibre aktivasyon nicelendirme (\S\ref{sec:method:q15}).
\emph{Uyuşma}: PyTorch-FP32 ile NumPy-Q15 argmax eşleşme oranı
(C-eşdeğer harness, bare-metal MCU değil; seed-0 MCU uyuşması
\%100, Tablo~\ref{tab:crossplatform}). Tohum~1, tüm aşamalarda
tutarlı yakınsama aykırısıdır (bkz. Şekil~\ref{fig:lowrank});
ortalama$\pm$SS bunu dahil eder. IHT maskesi $s{=}0.5$ seyrekliğini
zorunlu kıldığından tüm tohumlarda sıfır olmayan sayı ve bayt
aynıdır.}
\label{tab:multiseed}
\resizebox{\columnwidth}{!}{%
\begin{tabular}{@{}crrrrrr@{}}
\toprule
Tohum & LR F1 & Seyrek F1 & Q15/LUT F1 & Sıfır-olmayan & Bayt & Uyuşma \\
\midrule
0 & 0.918 & 0.921 & 0.913 & 283 & 566 & \%99.97 \\
1 & 0.781 & 0.680 & 0.663 & 283 & 566 & \%99.91 \\
2 & 0.901 & 0.893 & 0.895 & 283 & 566 & \%99.97 \\
3 & 0.890 & 0.895 & 0.891 & 283 & 566 & \%100.0 \\
4 & 0.904 & 0.890 & 0.904 & 283 & 566 & \%99.97 \\
\midrule
Ort$\pm$SS
  & $0.879{\pm}0.052$
  & $0.856{\pm}0.099$
  & $0.853{\pm}0.107$
  & 283 & 566 & --- \\
\bottomrule
\end{tabular}}
\end{table}

Parametre bütçesini karşılaştırmalı görmek için Tablo~\ref{tab:baselines},
gömülü hedefe aktarılan modeli aynı \hapt{} bölmesinde eğitilmiş bir MLP
taban çizgisiyle (en hafif yinelemeli olmayan referans) ve aynı
gizli boyuttaki LSTM ile GRU hücrelerinin kuramsal parametre
sayılarıyla karşılaştırır. Gömülü \fastgrnn{} LSQ modeli,
MLP taban çizgisinden \textbf{44$\times$ daha küçük} ve eşleşen
$H$ değerinde, sınıflandırıcı ek yükünden önce, bir LSTM hücresinden
\textbf{$\sim$5$\times$ daha küçüktür}.

\begin{table}[t]
\centering
\caption{Gömülü hedefe aktarılan modelin, yinelemeli olmayan ve yinelemeli
taban çizgilerine göre parametre ayak izi. LSTM ve GRU sayıları
$H{=}16$, $d{=}3$ için kuramsaldır (yalnızca hücre); ilişkili HAR
görevlerinde raporlanan doğrulukları \fastgrnn{} ile benzer ya da
biraz daha yüksektir, ancak $\sim$4$\times$ parametre dezavantajı
taşırlar~\cite{kusupati2018fastgrnn}. MLP F1 bu çalışmada
ölçülmüştür; sınıflandırıcı başlığı her durumda 102 parametre ekler.
\fastgrnn{} için: 181 hücre $+$ 102 başlık $=$ 283 toplam
sıfır-olmayan $\times$ 2~B $=$ 566~B.}
\label{tab:baselines}
\begin{tabular}{lrr}
\toprule
Model & Param. (hücre) & \hapt{} F1 \\
\midrule
MLP taban çizgisi              & 12{,}518 (tam)  & 0.847 \\
LSTM ($H{=}16$, kuramsal)      & 1{,}280          & --- \\
GRU ($H{=}16$, kuramsal)       & 960              & --- \\
\fastgrnn{} hücresi ($H{=}16$) & 338              & --- \\
\fastgrnn{} L (düşük rank)     & 328              & 0.879 \\
\fastgrnn{} LSQ (gömülü)          & \textbf{181}  & \textbf{0.918} \\
\bottomrule
\end{tabular}
\end{table}

\subsection{Çok Tohumlu Kararlılık ve Rank Seçimi}
\label{sec:results:lowrank}

Şekil~\ref{fig:lowrank}, yinelemeli rank değerleri
$r_u \in \{4, 6, 8, 12\}$ için her biri beş tohumdan oluşan test-F1
dağılımını raporlar. $r_u{=}8$ en iyi ortalamaya
($0.879 \pm 0.056$) ve en dar kuyruğa ulaşır. Dört yapılandırmanın
tamamında seed~$1$, ranktan bağımsız olarak tutarlı biçimde düşük
bir aykırı değerdir ($\sim$0.71~F1). Bu, özellikle vurgulanması
gereken bir yeniden üretilebilirlik bulgusudur: tek bir şanssız
başlatma, küçük örneklemli bir değerlendirmeyi yanıltabilir. Bu
model boyutunda en az beş tohum üzerinde ortalama~$\pm$~standart
sapma raporlamak bu nedenle isteğe bağlı değildir.

\begin{figure}[t]
\centering
\includegraphics[width=\linewidth]{figures/lowrank_seeds.pdf}
\caption{Yinelemeli rank seçimleri boyunca tohum-başına test F1.
$r_u{=}8$ en iyi ortalamaya ve kabul edilebilir varyansa ulaşır.
Kırmızı elmaslar rank-başına ortalamaları; yeşil bant seçilen
yapılandırmayı gösterir.}
\label{fig:lowrank}
\end{figure}

\subsection{Seyreklik Taraması}
\label{sec:results:sparsity}

Şekil~\ref{fig:sparsity}, sabit $r_u{=}8$ altında hedef seyreklik
$s \in \{0.3, 0.5, 0.7, 0.9\}$ değerlerini tarar. Tek tohumlu
tarama sığ ancak görünür bir U-eğrisi çizer: $s{=}0.5$ en fazla
doğruluğu korur, $s{=}0.3$ eğitim maliyetini haklı çıkaracak kadar
sıkıştırmaz ve $s \geq 0.7$ belirgin bozulmaya yol açar. Gömülü
uygulama için en uygun nokta olan $s{=}0.5$ değerinde beş tohumlu ortalama
$0.856 \pm 0.099$'dur; en yüksek tohum 0.92'ye ulaşır.
$s{=}0.5$ değerindeki yüksek standart sapma, seyrekleştirilmemiş
ortalamaya ($0.879 \pm 0.056$) göre bilgi-kuramsal sınırına yakın
bir parametreleştirme ile çalışmanın bedelidir.

\begin{figure}[t]
\centering
\includegraphics[width=\linewidth]{figures/sparsity_curve.pdf}
\caption{Hedef seyreklik karşısında test F1. Tek tohumlu tarama
(mavi çizgi) ve $s{=}0.5$ için beş tohumlu ortalama $\pm$ standart
sapma (kırmızı elmas). 0.4--0.6 bandı gömülü uygulama için en
uygun aralıktır.}
\label{fig:sparsity}
\end{figure}

\subsection{Nicelendirme Kipleri}
\label{sec:results:quant}

Şekil~\ref{fig:quant} ve Tablo~\ref{tab:quant}, seed-0 modelinde
dört nicelendirme yapılandırmasını karşılaştırır. Yalnızca ağırlık
Q15 kayıpsızdır: $\sigma$ ve $\tanh$ FP32'de çalışmayı sürdürür ve
ağırlık nicelendirme gürültüsü, eğitim yordamının kendi tohumlar
arası varyansının altında kalır. Kalibrasyonsuz saf Q15 aktivasyon
nicelendirmesi yıkıcıdır: F1 0.918'den 0.16'ya çöker ve
\textsc{standing} sınıfı tamamen kaybolur. Beş eğitim
mini-yığınından oluşan deterministik bir ön geçiş olan tensör-başına
aktivasyon kalibrasyonuyla Q15 çıkarımı, FP32 referans doğruluğunu
her sınıfta yuvarlama gürültüsü düzeyinde geri kazanır. Her iki
MCU'ya aktarılan yapılandırma budur.

\begin{figure}[t]
\centering
\includegraphics[width=\linewidth]{figures/quant_modes.pdf}
\caption{Seed-0 modelinde nicelendirme kipleri. Yalnızca ağırlık
Q15 kayıpsızdır; saf Q15 aktivasyonları modeli çökertir; kalibre
edilmiş Q15 aktivasyonları referansı yuvarlama gürültüsü içinde
geri kazanır.}
\label{fig:quant}
\end{figure}

\begin{table}[t]
\centering
\caption{Nicelendirme kipinin seed-0 test F1 üzerindeki etkisi (PyTorch
simülatöründen, Şek.~\ref{fig:quant} ile aynı rakamlar). Konuşlandırılan
C motoru ikinci satıra karşılık gelir: Q15 ağırlıklar, FP32
aktivasyonlar ve $\sigma$/$\tanh$ için 256 girişli LUT. Aynı 3{,}399
test penceresinde ayrıca ölçülen C-tarafı makro F1
\textbf{0.918}'dir; makalenin başlığında kullanılan değer budur.
Küçük PyTorch-C farkı, LUT'un aktivasyon hassasiyetini daha iyi
korumasından kaynaklanır. Dördüncü satır ise değerlendirilen ancak
konuşlandırılmayan tam-Q15 karşı-olgusudur.}
\label{tab:quant}
\setlength{\tabcolsep}{4pt}
\begin{tabular}{@{}lrl@{}}
\toprule
Kip & F1 & Rol \\
\midrule
Float32                                 & 0.914          & Referans \\
Q15 ağırlık, FP32 akt. (LUT)            & \textbf{0.914} & \textbf{Konuşlandırılan} (C: 0.918) \\
Q15 ağırlık, \emph{saf} Q15 akt.        & 0.16           & STANDING çöküşü \\
Q15 ağırlık, \emph{kalibre} Q15 akt.    & 0.913          & Karşı-olgu \\
\bottomrule
\end{tabular}
\end{table}

\subsection{Sınıf-Başına Davranış}
\label{sec:results:perclass}

Şekil~\ref{fig:perclass}, üç FP32 boru hattı aşaması ve
gömülü Q15 model boyunca sınıf-başına F1 değerlerini izler.
Statik sınıflar (\textsc{sitting}, \textsc{standing},
\textsc{laying}) yüksek F1 değerlerini korur. \textsc{walking} ve
\textsc{upstairs}, seyrekleştirme altında 1--3 yüzde puanı bozulur;
ancak 0.88'in üzerinde kalır. \textsc{downstairs}, tüm süreç boyunca
performansı belirleyen en zor sınıftır: düşük rank FP32'de $\sim$0.91'den seyrek
modelde 0.89'a düşer ve kalibre edilmiş Q15 sonrasında 0.91'e geri
döner. Bu bulgu, akıllı telefon düzeyi ivmeölçerlerde dinamik
iniş hareketini sistematik olarak en zor sınıf olarak tanımlayan
geniş HAR literatürüyle tutarlıdır~\cite{wang2019harsurvey,demrozi2020harsurvey}.

\begin{figure}[t]
\centering
\includegraphics[width=\linewidth]{figures/per_class_f1.pdf}
\caption{Sıkıştırma aşamaları boyunca sınıf-başına F1 (seed~0).
\textsc{laying} kolayca ayrılabilir; \textsc{downstairs} ise geniş
literatürü yansıtarak süreç boyunca en zor sınıf olarak kalır.}
\label{fig:perclass}
\end{figure}

\subsection{Platformlar Arası Deterministik Çıkarım}
\label{sec:results:crossplatform}

Aynı taşınabilir C kaynak dosyası, hem 8-bit \arduino{} hem de
16-bit \msp{} üzerinde değiştirilmeden derlenir. 3{,}399 pencerelik
test kümesinin tamamında tahminler olayların \%100'ünde uyuşur;
logitler ise mutlak olarak $10^{-2}$ değerinden daha iyi düzeyde
uyuşur (herhangi bir karar sınırı için anlamlı farkın oldukça
altında). Tablo~\ref{tab:hstate}, tek bir akış test penceresinde
altı zaman noktasında gizli bileşen $h_0$ değerini örnekler:
iki platform iki ondalık basamağa kadar aynı yörüngeleri üretir.
Bu açık bir sonuç değildir. İki aygıt farklı yazılım öykünmeli
kayan nokta uygulamaları, farklı çağrı kuralları ve farklı sözcük
boyutları kullanır. Buna rağmen uçtan uca tahminin bit-eşdeğer
olması, kalibre edilmiş Q15 ağırlıklar ile FP32'de biriktirip
sonra doyuran aritmetiğimizin uygulama ayrıntılarına karşı kararlı
olduğunu gösterir. Bu, güvenlikle ilişkili HAR için korunmaya değer
bir yeniden üretilebilirlik özelliğidir.

\begin{table}[t]
\centering
\caption{50~Hz zamanlamalı tek bir akış test penceresinde
platformlar boyunca gizli durum bileşeni $h_0$ evrimi.}
\label{tab:hstate}
\begin{tabular}{rrr}
\toprule
$t$ (örnek) & Arduino $h_0$ & MSP430 $h_0$ \\
\midrule
25  & $-0.720$  & $-0.720$  \\
50  & $-0.352$  & $-0.352$  \\
75  & $+0.459$  & $+0.459$  \\
100 & $+3.803$  & $+3.803$  \\
125 & $+11.388$ & $+11.388$ \\
128 & $+12.542$ & $+12.542$ \\
\bottomrule
\end{tabular}
\end{table}

\subsection{Gerçek Zamanlı Akış Performansı}
\label{sec:results:streaming}

Tablo~\ref{tab:streaming}, 50~Hz zamanlamalı akış altında
(20~ms bütçe) örnek-başına gecikmeyi raporlar. Her iki platform da
128 örnekli test penceresinin tamamında bütçeyi aşan örnek olmadan
gerçek zamanlı çalışır. Arduino ortalamada bütçenin \%46'sını
kullanır ve I\textsuperscript{2}C okumaları, LED geri bildirimi ve
UART günlüklemesi için geniş alan bırakır. \msp{} ortalamada
\%65 kullanır ve yine de bütçenin rahatça altında kalır.

Bölüm~\ref{sec:method:lut}'deki 256 girişli sigmoid/tanh LUT,
\msp{}'yu uygulanabilir kılan unsurdur. LUT olmadan \msp{} üzerinde
tam pencere çıkarımı $\sim$54~s sürer; LUT ile bu süre $\sim$1.8~s'ye
iner ve \textbf{30.5$\times$} hızlanma sağlar. Akış rejimine
çevrildiğinde bu, tamamen gerçek zamanlı olmayan bir çalışma ile
örnek başına yedi milisaniyelik zaman payı arasındaki
farktır. Aynı LUT \arduino{} üzerinde de etkindir; AVR çekirdeğinde
zaten donanım çarpıcısı bulunduğu için buradaki hızlanma daha
sınırlıdır (1.51$\times$).

\begin{table}[t]
\centering
\caption{50~Hz zamanlamalı akış performansı. Bütçe
20~ms/örnektir. Her iki platformda da 256 girişli LUT etkindir.}
\label{tab:streaming}
\begin{tabular}{lrrrr}
\toprule
Platform & Ort. & Maks. & Bütçe & Aşım \\
\midrule
\arduino{} & 9.21 & 10 & 46\% & 0 / 128 \\
\msp{}     & 13.0 & 14 & 65\% & 0 / 128 \\
\bottomrule
\end{tabular}
\end{table}

\begin{figure}[t]
\centering
\includegraphics[width=\linewidth]{figures/deploy_latency.pdf}
\caption{Üç referans platformda örnek-başına çıkarım gecikmesi.
Kırmızı kesikli çizgi, 50~Hz örneklemenin getirdiği 20~ms bütçeyi
gösterir. Her iki MCU hedefi de bütçenin belirgin biçimde altında
çalışır.}
\label{fig:latency}
\end{figure}

\subsection{Enerji Tüketimi}
\label{sec:results:energy}

Gecikme tablosunu tamamlamak için, elektriksel enerjiyi birincil
konuşlandırma hedefimiz ve değerlendirmedeki en kısıtlı kart olan
\msp{} üzerinde karakterize ediyoruz. Özel bir firmware modu
(\texttt{TEST\_MODE 3}), kartı sessiz bir ölçüm döngüsüne sokar ---
UART yok, LED yok, I\textsuperscript{2}C yok --- ve üç iş yükü
altında çalıştırır: \textit{Boşta} (düşük güç uykusu; geliştirme
kartının taban akımını yakalar), \textit{Akış50Hz} (her 20~ms'de bir
\texttt{fastgrnn\_step()} çağrısı; periyodun geri kalanı boşta;
konuşlandırılan HAR boru hattıyla aynı pacing) ve \textit{Sürekli}
(çıkarım üzerinde sıkı döngülenen en kötü durum zarfı). Akım, INA226
yüksek-taraf akım sensörüyle ölçüldü; aşağıdaki her değer, 60~s'lik
kararlı durum ortalamasıdır. Ayrıntılı protokol
\texttt{docs/energy\_measurement.md} dosyasında belgelenmiştir.

\begin{table}[t]
\centering
\caption{\msp{} üzerinde sessiz firmware ölçümlerinde gözlenen besleme
akımı (\texttt{TEST\_MODE~3}). VCC jumper'ı çıkarıldıktan sonra
LaunchPad hedef VCC rayı INA226 şöntünden (R\,=\,0{,}1\,$\Omega$,
adres~0x44) beslendiği için USB/debug köprüsü ölçülen rayın dışında
kalır. Değerler 60\,s sonrasındaki kararlı durum ortalamasıdır;
$V_{\text{cc}}$ yük altında ölçülen ray voltajıdır. Boşta LPM3
uykusudur (Timer\_A çalışır); $<$0{,}025\,mA girişi INA226 çözünürlük
tabanının altındadır (bu şönt için LSB\,=\,25\,$\mu$A).}
\label{tab:energy}
\begin{tabular}{lrrrr}
\toprule
Derleme & $V_{\text{cc}}$ & $I_{\text{boşta}}$ & $I_{50\text{Hz}}$ & $I_{\text{sürekli}}$ \\
        & (V)             & (mA)               & (mA)              & (mA) \\
\midrule
\msp{}, LUT    & 3{,}478 & $<$0{,}025 & 5{,}14 & 5{,}10 \\
\msp{}, LUT yok & 3{,}478 & $<$0{,}025 & \multicolumn{1}{c}{---} & 5{,}08 \\
\bottomrule
\end{tabular}
\end{table}

\begin{table}[t]
\centering
\caption{\msp{} konuşlandırması için türetilmiş enerji.
$E$/çıkarım\,=\,$P_{\text{sürekli}} \times t_{\text{adım}}$;
$E$/pencere\,=\,$P \times T_{\text{pencere}}$; 50\,Hz akış satırında
LPM uykusu araya girer (128 örnek $\times$ 20\,ms\,=\,2{,}56\,s).
Pil ömrü 2000\,mAh / 3{,}7\,V Li-ion hücre (7{,}4\,Wh) varsayar.
LUT'suz satır ablasyon amaçlıdır: LUT olmadan 20\,ms akış zaman
sınırı aşılır ($t_{\text{adım}}\approx421$\,ms), dolayısıyla 50\,Hz
akış satırı yoktur.}
\label{tab:energy-derived}
\resizebox{\columnwidth}{!}{%
\begin{tabular}{llrrr}
\toprule
Derleme & Mod & $E$ / çıkarım & $E$ / pencere & Pil ömrü \\
        &     & ($\mu$J)      & (mJ)          & (s) \\
\midrule
\msp{}, LUT    & 50\,Hz akış & 246     & 31{,}5 & 602 \\
\msp{}, LUT    & sürekli     & 246     & ---    & 417 \\
\msp{}, LUT yok & sürekli   & 7{,}440 & ---    & 419 \\
\bottomrule
\end{tabular}}
\end{table}

Bu enerji rakamları, Bölüm~\ref{sec:intro}'deki motivasyonu
tamamlamaktadır: 566~bayt Flash'a sığan ve 50~Hz'de gerçek zamanlı
çalışan aynı model, donanım çarpıcısı olmayan \msp{} üzerinde,
düğme pille çalışan giyilebilir bir cihazın bütçesinin de oldukça
altında bir akım çekmektedir. Aynı örnekleme hızında buluta akışlı
bir alternatif ise --- veri hızıyla artan görev döngüleriyle ---
radyo aktif iletim akımına ($\sim$10--20~mA) tâbi olur ve ölçtüğümüz
her eksende yerel çıkarım yoluna göre dezavantajlı kalır. LUT'suz
sürekli satırı ayrıca aktivasyon tablosunun etkisini izole eder:
LUT olmadan aynı firmware 50~Hz zamanlamasını karşılayamaz, bu yüzden
enerjisi konuşlandırılabilir bir akış modu değil, bir ablasyon olarak
raporlanır.
